%% EXEMPLO FICTÍCIO, Atividade 8: publicação, que conta à parte
\begin{atividade}{Artigo completo sobre a fonte de corrente do estimulador, em congresso nacional}{
  tipo           = pesquisa,
  aproveitamento = PUB,
  publicacao     = nacional,
  horas          = 40,
  inicio         = 2024-11,
  fim            = 2025-05,
  projeto        = {NEBULA, estimulação elétrica funcional para ciclismo},
  papel          = {Primeira autora},
  supervisor     = {Prof. Dr. João Modelo, coautor e orientador do trabalho},
  registros      = {NBL-161; NBL-168},
  competencias   = {II, V, VIII},
  disciplinas    = {Eletrônica Analógica; Instrumentação Eletrônica},
}

\begin{contexto}
Os resultados da revisão B da placa de potência (Atividades~1, 2 e~4) resolviam um problema que outras
equipes de estimulação elétrica funcional também têm: manter a corrente quando a impedância do eletrodo
sobe ao longo de uma sessão. O supervisor do projeto e o professor sugeriram transformar o projeto e os
ensaios num artigo, para que a solução ficasse publicada e citável.

Esta atividade é a escrita e a publicação do artigo. O projeto, a simulação e os ensaios que ele relata
são as Atividades~1, 2 e~4, que contam na participação; aqui contam só as horas de escrever, revisar e
apresentar (art.~7\textordmasculine\ da resolução).
\end{contexto}

\begin{contribuicao}
Sou a primeira autora. Escrevi a primeira versão inteira, refiz as figuras a partir dos dados brutos
dos ensaios, respondi aos dois pareceristas e apresentei o trabalho na sessão oral. O professor revisou
o texto e a discussão; o supervisor do projeto e um colega do subsistema, coautores, conferiram os
números e as figuras.
\end{contribuicao}

\begin{desenvolvimento}
\fichatecnica{
  pergunta = {Uma fonte de corrente com amplificador operacional e MOSFET, alimentada a \SI{120}{\volt}, mantém pulsos de \SI{100}{\milli\ampere} com eletrodos degradados?},
  metodo   = {Projeto, simulação do transitório com o modelo de Randles e ensaio em bancada com carga equivalente e com eletrodos usados},
  bases    = {IEEE Xplore e PubMed, com os termos functional electrical stimulation, current source e electrode impedance},
  produto  = {Artigo completo, oito páginas, nos anais de um congresso nacional de engenharia biomédica (fictício)},
  situacao = {Publicado; DOI \cod{10.0000/exemplo.2025.0042}},
}

O artigo segue a estrutura do relatório de teste da placa (\cod{NRO-PRO-003-901}), reescrita para quem
não conhece o projeto: a introdução situa o problema na literatura \parencite{merrill2005}, a seção de
métodos descreve o circuito e o modelo da carga, e a de resultados compara a simulação com as medições,
com a incerteza de cada uma. A revisão por pares pediu duas mudanças: a análise do pior caso de
dissipação, que acrescentei com os números da Atividade~1, e a declaração de que nenhum dado de pessoa
foi usado, já que os ensaios com o atleta-piloto ficaram de fora do artigo.
\end{desenvolvimento}

\begin{resultados}
O artigo foi aceito na primeira rodada, com revisões menores, e apresentado em sessão oral de quinze
minutos. É a primeira publicação da equipe sobre o hardware do estimulador, e o repositório da placa
passou a citá-lo. Como publicação em evento nacional, rende 1 crédito (art.~9\textordmasculine, II), à
parte da participação.
\end{resultados}

\begin{evidencias}
\evidencia{Publicação}{DOI 10.0000/exemplo.2025.0042}{O artigo completo, com a lista de autores}{Anais do congresso (exemplo)}
\evidencia{Comprovante}{Certificado de apresentação (exemplo)}{A apresentação oral do trabalho}{Anexo do relatório}
\evidencia{Atividade}{NBL-161}{O cartão da escrita do artigo, atribuído a mim e concluído}{SOMA, Atividades}
\end{evidencias}

\begin{aprendizados}
Escrever para quem não estava no laboratório me obrigou a justificar o que, no relatório interno, era
óbvio para a equipe. A resposta aos pareceristas ensinou a separar o que o ensaio mostra do que eu
gostaria que ele mostrasse. As disciplinas de eletrônica deram o conteúdo; a de metodologia, que eu
tinha achado abstrata, ganhou sentido aqui.
\end{aprendizados}
\end{atividade}
